\subsection{A natural obstruction on three-dimensional concordance components}

The same relative class defines a homomorphism on all ordinary
concordance components, without an orientability assumption.
We use the real, joint-family, and free-top conventions of
Theorem~\ref{thm:cyclic-three-concordance-obstruction}.

\begin{theorem}\label{thm:natural-three-concordance-obstruction}
For every connected closed real PL three-manifold $B$ there
is a natural group homomorphism
\begin{equation}\label{eq:natural-three-concordance-class}
 c_B:\pi_0P_{\mathrm{TOP}}(B)\longrightarrow
                         H^2(B;\mathbb Z/2)
\end{equation}
which vanishes on the PL image and on the definable image
for every fixed o-minimal expansion of $\mathbb R$.
For a PL homeomorphism $f:B\to B'$ its naturality is
\begin{equation}\label{eq:three-concordance-typed-naturality}
 c_{B'}\bigl([(f\times1)H(f^{-1}\times1)]\bigr)
                    =(f^{-1})^*c_B([H]).
\end{equation}
If $B$ is oriented and $\pi_1(B)$ is finite cyclic, then
$c_B$ is surjective.  In particular, for even cyclic order,
\[
 \pi_0P_{\mathrm{TOP}}(B)/\ker c_B\cong\mathbb Z/2,
 \qquad
 \operatorname{im}\pi_0P_{\mathrm{def}}(B)\subseteq\ker c_B.
\]
\end{theorem}

\begin{proof}
Put $X=B\times I$, and abbreviate mod-two cohomology by
$H^*(-)$.  Let $R_B$ be the subgroup of TOP concordances
whose top map is PL.  For $h\in R_B$, both boundary maps of
$X$ are PL, so the mapping cylinder of $h$ has prescribed
PL tangent reductions on its entire boundary.  Define
\begin{equation}\label{eq:three-boundary-pl-class}
 k_B(h)=\sigma^{-1}\operatorname{ks}(M_h,\partial M_h)
                     \in H^3(X,\partial X).
\end{equation}
This extends \eqref{eq:relative-three-base-cs}; the side of
$M_h$ is now the mapping cylinder of the PL boundary map.

\emph{Common product collars.}
The boundary identification can be made literal on a product
collar.  If $a$ is the top map of $h$ and $0<r<1/2$, let
$e_r(x,s)=(x,r+(1-2r)s)$ and define
\[
 h_r(x,s)=
 \begin{cases}
  (x,s),&0\leq s\leq r,\\
  e_rhe_r^{-1}(x,s),&r\leq s\leq1-r,\\
  (a(x),s),&1-r\leq s\leq1.
 \end{cases}
\]
The pieces agree at both interfaces, and the inverse is the
same construction applied to $h^{-1}$.  Setting $h_0=h$
gives an ordinary isotopy with fixed boundary maps.  In a
product metric, uniform continuity bounds the displacement
from $h$ by $r+\omega_h(r)$; the same estimate holds for
the inverse.  A compact parameter family has a common
modulus of continuity.  Thus the construction works for
families as well.  At a fixed width it satisfies
$(uv)_r=u_rv_r$: inside it is conjugation by $e_r$, and
outside it is multiplication of the boundary maps.  A fixed
width also preserves joint PL families.  These collars let
us compare or compose any finite collection with identical
normal-line conventions.

\emph{The relative difference rule.}
The standard PL structure of $X$ supplies a section of the
stable tangent-reduction fibration.  Its fiber is
\[
 \mathrm{TOP}/\mathrm{PL}\simeq K(\mathbb Z/2,3)
\]
by \cite[Essay~V, Theorem~5.5(I), p.~251]{KirbySiebenmann1977}.
This coefficient system is untwisted, even when $B$ is
nonorientable.  Indeed the chosen section points each fiber;
the pointed self-equivalence space of $K(\mathbb Z/2,3)$
has one component, since $\operatorname{Aut}(\mathbb Z/2)=1$,
and its positive homotopy groups are
\[
 \widetilde H^{3-i}(K(\mathbb Z/2,3);\mathbb Z/2)=0
 \qquad(i>0).
\]
Consequently relative differences of its sections form the
ordinary group $H^3(X,\partial X)$.

Pull back the standard reduction by $h$ using the canonical
tangent-microbundle identification, with its prescribed PL
boundary identification.  Its difference from the original
reduction is \eqref{eq:three-boundary-pl-class}.  Inserting
an intermediate reduction in the comparison for a composite
gives
\begin{equation}\label{eq:three-relative-cocycle}
 k_B(uv)=k_B(v)+v^*k_B(u),\qquad u,v\in R_B.
\end{equation}
The boundary identifications compose because they are PL.
This is the relative PL difference-of-reductions argument
underlying the classical Casson--Sullivan cocycle; compare
\cite[Proposition~2.23 and its full proof, pp.~2508--2509]{Galvin2026CassonSullivan}.
It also proves invariance under an ordinary isotopy with
fixed PL boundary maps, and under one whose boundary is a
jointly PL family.  The common collars above provide the
literal relative boundary convention for such a family.
If $h$ is PL, its mapping cylinder is PL, so $k_B(h)=0$.

Every bottom-fixed concordance acts trivially on
$H^3(X,\partial X)$.  To see this without orientability,
its top map $a$ is homotopic to the identity by
$(x,t)\mapsto\operatorname{pr}_B h(x,t)$.  The restriction
maps $H^j(X)\to H^j(\partial X)$ are diagonal for $j=2,3$,
and the latter diagonal is injective.  The exact sequence
of the pair therefore gives
\begin{equation}\label{eq:three-relative-trivial-action}
 H^3(X,\partial X)
 \cong\operatorname{coker}\bigl(H^2(B)
                 \xrightarrow{\Delta}H^2(B)\oplus H^2(B)\bigr).
\end{equation}
The boundary action is the identity on the bottom and $a^*=1$
on the top.  It is therefore the identity on this cokernel.
Thus \eqref{eq:three-relative-cocycle} is additive on $R_B$.

\emph{Endpoint traces have zero class.}
If $\delta_s$ is an ordinary path from the identity to a PL
map, set $T_\delta(x,s)=(\delta_s(x),s)$.  The
three-manifold automorphism comparison
\eqref{eq:three-manifold-full-automorphism-comparison}
replaces $\delta$ by a jointly PL path with the same
endpoints and the same ordinary relative path class, as in
the preceding proof.  Its trace is PL, and the intervening
trace isotopy fixes both boundary maps.  Hence
\begin{equation}\label{eq:three-trace-class-zero}
 k_B(T_\delta)=0.
\end{equation}

\emph{Normalization and independence.}
Let $h\in P_{\mathrm{TOP}}(B)$ have top map $f$.  Surjectivity
on components of the automorphism comparison supplies a
PL map $a$ and an ordinary path $\gamma_s:f\to a$.  Put
\begin{equation}\label{eq:three-endpoint-normalization}
 E_\gamma(x,s)=(\gamma_s f^{-1}(x),s),
 \qquad h^\gamma=E_\gamma h\in R_B.
\end{equation}
Replacing $\gamma_s$ by $\gamma_{ts}$ gives a path from the
identity to $E_\gamma$ in the bottom-fixed group.  Thus
$h^\gamma$ belongs to the same ordinary free-top component
as $h$.  If $\gamma':f\to a'$ is another choice, then
\[
 h^{\gamma'}(h^\gamma)^{-1}=E_{\gamma'}E_\gamma^{-1}
            =T_\delta,\qquad \delta_s=\gamma'_s\gamma_s^{-1}.
\]
This path starts at the identity and ends at the PL map
$a'a^{-1}$.  Additivity and \eqref{eq:three-trace-class-zero}
give $k_B(h^{\gamma'})=k_B(h^\gamma)$.

\emph{Ordinary component invariance.}
Let $H_t$ be a path with $H_0=h'\in R_B$, $H_1=k$, and
top path $f_t$ starting at the PL map $a=f_0$.  Define
\[
 T_t(x,s)=(f_{ts}a^{-1}(x),s),\qquad Q_t=T_t^{-1}H_t.
\]
The family $Q_t$ fixes the bottom and has the constant PL
top map $a$.  Its relative class is therefore constant.
At $t=1$ it is the normalization of $k$ obtained from
\[
 \gamma_s=a f_s^{-1}f_1:f_1\longrightarrow a,
 \qquad E_\gamma(x,s)=(a f_s^{-1}(x),s)=T_1^{-1}(x,s).
\]
Hence the normalized value is constant along $H_t$.  For a
path whose initial map is not normalized, first prepend
the reverse of its normalization path.  This proves that
the value is well-defined on all ordinary components.

Compose it with the interval isomorphism
$H^3(X,\partial X)\cong H^2(B)$ to define $c_B$.
For arbitrary $h,k$, multiply their two normalization
paths.  The product of their normalized representatives
lies in $R_B$ and represents the component of $hk$.
Additivity on $R_B$ yields
\[
 c_B([hk])=c_B([h])+c_B([k]).
\]
For a PL concordance choose a constant normalization path,
giving zero.  Theorem~\ref{result-4} makes every definable
concordance ordinarily isotopic to a PL one, so its class
also vanishes.

Conjugation by $f\times1$, for a PL homeomorphism of bases
$f:B\to B'$, transports the product reductions, boundary
markings, and mapping cylinders.  Pullback of the relative
obstruction through the inverse change of coordinates is
$(f^{-1}\times1)^*$.  Naturality of interval suspension
then gives exactly \eqref{eq:three-concordance-typed-naturality}.

Finally, if $B$ is oriented with finite cyclic fundamental
group, Lemma~\ref{lem:compact-cyclic-relative-cs-realization}
realizes every class by a map fixing both ends.  Such a map
is already normalized, so $c_B$ is surjective.  For even
cyclic order, $H^2(B;\mathbb Z/2)=\mathbb Z/2$ by mod-two
duality and abelianization.  The first isomorphism theorem
gives the asserted quotient by $\ker c_B$.
\end{proof}

The cocycle and its vanishing on PL maps are classical
formal Casson--Sullivan properties.  The construction above
records their descent to the specified ordinary free-top
component group. Determining which zero-class components contain
a PL representative remains a separate question.
